\documentclass[10pt]{beamer}

% Sprache
\usepackage[ngerman]{babel}

% Custom Theme (liegt im selben Verzeichnis)
\usetheme{MLUProdLog}

\setbeamersize{text margin left=0.5cm,text margin right=0.5cm}

% Pakete
\usepackage[linesnumbered,plainruled,nofillcomment]{algorithm2e}
\usepackage{amsmath}
\usepackage{amsfonts}
\usepackage{amssymb}
\usepackage{amsthm}
\usepackage{array}
\usepackage{bm}
\usepackage{bbm}
\usepackage{enumitem}
\setitemize{label=\usebeamerfont*{itemize item}\usebeamercolor[fg]{itemize item}\usebeamertemplate{itemize item}}
\usepackage{eurosym}
\usepackage{rotating}
\usepackage[utf8]{inputenc}
\usepackage{natbib}
\usepackage{tabularx}
\usepackage[T1]{fontenc}
\usepackage{dcolumn}
\usepackage{booktabs}
\usepackage{hyperref}
\usepackage{ifthen}
\usepackage{colortbl}
\usepackage{multirow}
\usepackage{subfig}
\usepackage{lmodern}
\usepackage{courier}
\usepackage{siunitx}
\usepackage{tikz}
\usepackage{relsize}
\usepackage[overlay,absolute]{textpos}
\setlength{\TPHorizModule}{10mm}
\setlength{\TPVertModule}{\TPHorizModule}

\usetikzlibrary{fit,shapes.geometric}

\sisetup{
    locale=US,
    group-minimum-digits = 4,
    group-digits = integer,
    group-separator ={,},
    detect-weight=true,
    detect-inline-weight=math,
    output-decimal-marker = {.}
}

\robustify\bfseries
\usefonttheme[onlymath]{serif}

\def\newblock{\hskip .11em plus .33em minus .07em}

\newcommand\mc[1]{\multicolumn{1}{c}{#1}}
\DeclareMathOperator{\di}{d\!}

% =========================================================
% Fraunhofer IKTS Logo als TikZ-Makro
% Offizielle Fraunhofer-Farbe: #179C7D
% Verwendung: \FraunhoferIKTSlogo[Hoehe]
% =========================================================
\definecolor{FhGreen}{HTML}{179C7D}

\newcommand{\FraunhoferIKTSlogo}[1][8ex]{%
  \begin{tikzpicture}[baseline=(current bounding box.center)]
    % Skalierungsfaktor abhaengig von gewuenschter Hoehe
    \pgfmathsetmacro{\sc}{0.38}
    % Grünes Quadrat (Fraunhofer-Signet)
    \fill[FhGreen] (0,0) rectangle (\sc,\sc);
    % Weisser Streifen horizontal (oberes Drittel)
    \fill[white] (0,0.67*\sc) rectangle (\sc,\sc);
    % Weisser Streifen vertikal (rechtes Drittel)
    \fill[white] (0.67*\sc,0) rectangle (\sc,\sc);
    % Grünes Rechteck oben rechts (Fraunhofer-F)
    \fill[FhGreen] (0.67*\sc,0.67*\sc) rectangle (\sc,\sc);
    % Schriftzug "Fraunhofer" bold
    \node[anchor=west, font=\sffamily\bfseries\tiny, text=FhGreen]
      at (\sc+0.04,0.72*\sc) {Fraunhofer};
    % Schriftzug "IKTS" kleiner
    \node[anchor=west, font=\sffamily\small\bfseries, text=FhGreen]
      at (\sc+0.04,0.28*\sc) {IKTS};
  \end{tikzpicture}%
}

% Math-Befehle
\newcommand{\Y}{\mathbf{Y}}
\newcommand{\bZ}{\mathbf{Z}}
\newcommand{\X}{\mathbf{X}}
\newcommand{\V}{\mathbf{V}}
\newcommand{\Q}{\mathbf{Q}}
\newcommand{\A}{\mathbf{A}}
\newcommand{\B}{\mathbf{B}}
\newcommand{\Db}{\mathbf{D}}
\newcommand{\y}{\mathbf{y}}
\newcommand{\f}{\mathbf{f}}
\newcommand{\x}{\mathbf{x}}
\newcommand{\bh}{\mathbf{h}}
\newcommand{\pb}{\mathbf{p}}
\newcommand{\rr}{\mathbf{r}}
\newcommand{\z}{\mathbf{z}}
\newcommand{\Z}{\mathbf{Z}}
\newcommand{\lb}{\mathbf{l}}
\newcommand{\rb}{\mathbf{r}}
\newcommand{\bs}{\mathbf{s}}
\newcommand{\bt}{\mathbf{t}}
\newcommand{\bu}{\mathbf{u}}
\newcommand{\E}{\mathbf{E}}
\newcommand{\Sb}{\mathbf{S}}
\newcommand{\bd}{\mathbf{d}}
\newcommand{\e}{\mathbf{e}}
\newcommand{\ba}{\mathbf{a}}
\newcommand{\bb}{\mathbf{b}}
\newcommand{\bz}{\mathbf{z}}
\newcommand{\balpha}{\boldsymbol{\alpha}}
\newcommand{\bphi}{\boldsymbol{\phi}}
\newcommand{\bpi}{\boldsymbol{\pi}}
\newcommand{\bPi}{\boldsymbol{\Pi}}
\newcommand{\bkappa}{\boldsymbol{\kappa}}
\newcommand{\bnu}{\boldsymbol{\nu}}
\newcommand{\bomega}{\boldsymbol{\omega}}
\newcommand{\bOmega}{\boldsymbol{\Omega}}
\newcommand{\bXi}{\boldsymbol{\Xi}}
\newcommand{\bmu}{\boldsymbol{\mu}}
\newcommand{\bupsilon}{\boldsymbol{\upsilon}}
\newcommand{\bUpsilon}{\boldsymbol{\Upsilon}}
\newcommand{\bSigma}{\boldsymbol{\Sigma}}
\newcommand{\bpsi}{\boldsymbol{\psi}}
\newcommand{\btheta}{\boldsymbol{\theta}}
\newcommand{\bepsilon}{\boldsymbol{\epsilon}}
\newcommand{\bdelta}{\boldsymbol{\delta}}
\newcommand{\bbeta}{\boldsymbol{\beta}}
\newcommand{\bxi}{\boldsymbol{\xi}}
\newcommand{\bPhi}{\boldsymbol{\Phi}}
\newcommand{\bTheta}{\boldsymbol{\Theta}}
\newcommand{\bvarepsilon}{\boldsymbol{\varepsilon}}
\newcommand{\bgamma}{\boldsymbol{\gamma}}
\newcommand{\bGamma}{\boldsymbol{\Gamma}}
\newcommand{\blambda}{\boldsymbol{\lambda}}
\newcommand{\Rre}{\mathbb{R}}
\newcommand{\Nre}{\mathbb{N}}
\newcommand{\Nrm}{\mathrm{N}}
\newcommand{\Ac}{\mathcal{A}}
\newcommand{\Bc}{\mathcal{B}}
\newcommand{\bc}{\Beta}
\newcommand{\Cc}{\mathcal{C}}
\newcommand{\Dc}{\mathcal{D}}
\newcommand{\Mc}{\mathcal{M}}
\newcommand{\Ec}{\mathcal{E}}
\newcommand{\Fc}{\mathcal{F}}
\newcommand{\Gc}{\mathcal{G}}
\newcommand{\Erm}{\mathrm{E}}
\newcommand{\Lc}{\mathcal{L}}
\newcommand{\Nc}{\mathcal{N}}
\newcommand{\Nbc}{\mathcal{\bar{N}}}
\newcommand{\Oc}{\mathcal{O}}
\newcommand{\Pc}{\mathcal{P}}
\newcommand{\Rc}{\mathcal{R}}
\newcommand{\Sc}{\mathcal{S}}
\newcommand{\ssc}{\mathcal{s}}
\newcommand{\Hc}{\mathcal{H}}
\newcommand{\Ic}{\mathcal{I}}
\newcommand{\Kc}{\mathcal{K}}
\newcommand{\Tcc}{\mathcal{T}}
\newcommand{\Vc}{\mathcal{V}}
\newcommand{\Jc}{\mathcal{J}}
\newcommand{\Zc}{\mathcal{Z}}
\newcommand{\Yc}{\mathcal{Y}}
\newcommand{\Xc}{\mathcal{X}}

\newcommand{\format}[1]{\ensuremath{\mathrm{#1}}}
\newcommand{\p}{\phantom{1}}
\newcommand{\pp}{\phantom{11}}
\newcommand{\ppp}{\phantom{111}}
\newcommand{\pppp}{\phantom{1111}}
\newcommand{\sw}[1]{\rotatebox{45}{\parbox{1cm}{\baselineskip0cm#1}}}

\newcolumntype{C}[1]{>{\centering\arraybackslash}m{#1}}
\newcolumntype{L}[1]{>{\raggedright\arraybackslash}m{#1}}
\newcolumntype{R}[1]{>{\raggedleft\arraybackslash}m{#1}}
\newcolumntype{Y}{>{\centering\arraybackslash}X}

\newcommand{\CC}[2]{\multicolumn{#1}{@{}c@{}}{#2}}
\newcommand{\pc}{\tiny\%}

% =========================================================
% Titel-Informationen
% Logo: Fraunhofer IKTS TikZ-Makro auf allen Folien (Fusszeile/Header via \logo)
% =========================================================
\logo{\FraunhoferIKTSlogo[5ex]}
\title[Agroforst-Wertsch\"opfungsketten]{Planung von Agroforst-Wertsch\"opfungsketten}
\author[Thomas Kirschstein]{\textbf{Thomas Kirschstein}}
\date{}

\setbeamercovered{invisible}

% ************************************* DOCUMENT **************************************************
\begin{document}

% ----------------------------------------- TITELFOLIE -----------------------------------------------------
\frame[plain,label=title]{
  \titlepage
  % Fraunhofer IKTS Logo prominent auf der Titelfolie (unten rechts)
  \begin{textblock}{4}(8.5,7.5)
    \FraunhoferIKTSlogo[10ex]
  \end{textblock}
}

% ----------------------------------------- BLANKO-FOLIE -----------------------------------------------------
\begin{frame}{}
  % Blanko-Folie (frei belegbar)
  \vfill
  \vfill
\end{frame}

% ########################################  ABSCHNITTE #############################################

\section{Einleitung}

\begin{frame}{Motivation}
  \begin{itemize}
    \item Agroforst kombiniert Geh\"olze (B\"aume, Str\"aucher) mit landwirtschaftlichen Kulturen oder Vieh auf derselben Fl\"ache
    \item Vielf\"altige Produkte: Holz, Fr\"uchte, N\"usse, Biomasse, Nahrungsmittel, Futtermittel
    \item Herausforderung: Planung effizienter Wertsch\"opfungsketten unter Ber\"ucksichtigung langer Produktionszyklen
  \end{itemize}
\end{frame}

% ----------------------------------------- NEXT FRAME -----------------------------------------------------

\section{Wertsch\"opfungskette}

\begin{frame}{Struktur der Agroforst-Wertsch\"opfungskette}
  \begin{columns}[T]
    \begin{column}{.5\textwidth}
      \textbf{Produktionsseite:}
      \begin{itemize}
        \item Standortauswahl \& Fl\"achennutzung
        \item Baumarten- \& Kulturwahl
        \item Ernte- \& Pflegema\ss{}nahmen
      \end{itemize}
    \end{column}
    \begin{column}{.5\textwidth}
      \textbf{Logistik \& Vermarktung:}
      \begin{itemize}
        \item Transport \& Lagerung
        \item Verarbeitungsstufen
        \item Absatzkan\"ale (regional, \"uberregional)
      \end{itemize}
    \end{column}
  \end{columns}
\end{frame}

% ----------------------------------------- NEXT FRAME -----------------------------------------------------

\section{Optimierungsmodell}

\begin{frame}{MILP-Formulierung}
  \begin{block}{Notation}
    \begin{tabularx}{.99\textwidth}{l @{\hskip1ex} X}
      $\Fc$ & Menge der Fl\"achen \\
      $\Pc$ & Menge der Produkte \\
      $\Tcc$ & Planungshorizont \\
      $x_{f,p,t} \in \{0,1\}$ & 1, wenn Fl\"ache $f$ in Periode $t$ mit Produkt $p$ bebaut \\
      $c_{f,p,t}$ & Kosten der Bepflanzung \\
    \end{tabularx}
  \end{block}
  \vspace{2ex}
  \begin{equation*}
    \min \rightarrow \sum_{f \in \Fc} \sum_{p \in \Pc} \sum_{t \in \Tcc} x_{f,p,t} \cdot c_{f,p,t}
  \end{equation*}
\end{frame}

% ----------------------------------------- NEXT FRAME -----------------------------------------------------

\section{Zusammenfassung}

\begin{frame}{Fazit}
  \begin{itemize}\itemsep1ex
    \item Planung von Agroforst-Wertsch\"opfungsketten erfordert integrierte Optimierungsans\"atze
    \item Ber\"ucksichtigung langer Produktionszyklen und Multiprodukt-Strukturen
    \item[$\rightarrow$] MILP-Modelle erm\"oglichen systemweite Optimierung
  \end{itemize}
  \vspace{2ex}
  \textbf{Offene Forschungsfragen:}
  \begin{itemize}
    \item Skalierbarkeit bei gro\ss{}en Instanzen
    \item Integration von Unsicherheiten (Ertr\"age, Preise, Klima)
    \item Mehrzielige Optimierung (\"okonomisch \& \"okologisch)
  \end{itemize}
\end{frame}

% ----------------------------------------- SCHLUSSBEMERKUNG -----------------------------------------------------
\begin{frame}{Abschlie\ss{}end\ldots}
  \vspace{10ex}
  \vfill\LARGE
  Fragen?\\[5ex]
  \vfill
\end{frame}

%--------------------------------------------------------------------
\appendix
\begin{frame}{Literatur}
  \tiny
  \bibliographystyle{abbrvnat}
  \bibliography{literature}
\end{frame}

\end{document}
